\subsection{Sprint 3}

O dia da entrega da sprint passada contou com talvez metade do time presente,
e com uma entrega dita pelos AGES IV do projeto como a “\textit{pior entrega que já viram
em um projeto da AGES}”. Ao invés da retrospectiva, foi decidido realizar uma reunião
de emergência pelos presentes para conversar sobre o que estava acontecendo, o
que devia mudar e talvez como proceder.

Feita a reunião e todos tendo dito suas opiniões, o time resolveu abandonar
certas ideias do \textit{agile} visto que a rotina de trabalho não estava fluindo como esperado.
Os AGES IV tomaram uma medida mais direta na atribuição de duplas e tasks, tendo
maior controle sobre o projeto em si.

Continuei no \textit{frontend} esta sprint e de dupla com o AGES II Otávio Cunha.
Fomos encarregados de resolver algumas dívidas técnicas como responsividade na
tela de cadastro do usuário e alguns componentes, além de desenvolver novas telas
sobre visualização de documentos em procedimentos criminais.

Acabamos mexendo bastante especificamente em CSS para resolver as
dívidas de responsividade, já na parte de visualização de documentos conseguimos
trabalhar mais tranquilamente. O que mais nos tomou tempo foi integrar corretamente
as rotas do backend para chamar os documentos de cada procedimento.

Mesmo com estas mudanças a entrega ainda assim precisou de ajustes de
última hora dos AGES III e IV, devido a falta de testes pelo time na hora de integrar
as suas tarefas.
\\

\fbox{%
    \parbox{0.95\linewidth}{%
        \setlength{\parindent}{15pt}
            \itshape    A reunião do dia da entrega realmente foi quando o time virou a chave, após
                este dia todo pareciam mais envolvidos com o projeto tanto quanto nas reuniões
                diárias e nos canais de comunicação. Aqui  também, percebi bastante meu
                melhor desempenho e atitude.
                
                Todavia, mesmo com as melhorias aparentes a entrega constou com
                mudanças locais pelos AGES IV. Notei que alguns AGES I e II ainda pecavam na
                hora de testar cada uma de suas tarefas com o projeto todo.
    }%
}